% L15a strategy deck: the minimum-variance weighting rule.
% One short deck per rule on the strategy menu, read next to the rule's code file
% (strategies/weighting/minimum_variance.jl). Written by make_strategy_decks.py (session scratchpad, 2026-11-27) and then
% kept by hand: edit this file, not the generator. Build with make in this folder: LaTeX runs from strategies/,
% so paths are relative to it.
\documentclass[aspectratio=169,11pt]{beamer}
\usepackage{vnslides}

\newcommand{\repo}{https://github.com/varnerlab/CHEME-5660-CourseRepository-Fall-2026/blob/main/lectures/week-15/L15a}

\title{\texorpdfstring{{\vnheavy L15a}}{L15a}: Minimum Variance}
\subtitle{Weighting Rule, CHEME 5660 Strategy Menu, Cornell University}
\author{Copyright Jeffrey Varner 2026. All Rights Reserved}

\begin{document}

\vntitlepage

\begin{frame}{Disclaimer and Risks}
  \vspace{0.6em}
  \textbf{\ink{Training only.}} This content is for training and informational
  purposes only. The teaching team does not make, give, or endorse any offer or
  solicitation to buy or sell securities or derivative products, or provide
  investment or trading advice or strategy.

  \vspace{1.1em}
  \textbf{\ink{Risk.}} Trading involves risk and is generally not recommended for
  those with limited resources, experience, or a low risk tolerance. Past
  performance does not guarantee future success.

  \vspace{1.1em}
  \textbf{\ink{Responsibility.}} You are responsible for your investment decisions
  based on your financial circumstances, objectives, risk tolerance, and
  liquidity needs.
\end{frame}

\begin{frame}{The Idea}
  \vspace{0.3em}
  The minimum-variance rule holds the long-only portfolio of the basket with the smallest variance of its
  growth rate:
  \begin{itemize}\setlength{\itemsep}{0.5em}
    \item \hilite{Long only}: the rule never sells a ticker short.
    \item \hilite{Risk only}: the rule ignores the preference weights and the expected growth rates.
    \item \hilite{Covariance}: the single index model gives the covariance of the tickers' growth rates.
  \end{itemize}

  \vspace{0.5em}
  The basket comes from the picking rule.
\end{frame}

\begin{frame}{The Minimum-Variance Weights}
  \vspace{0.3em}
  Under the single index model, the growth rates of tickers $i$ and $j$ have the covariance
  \[
    \Sigma_{ij}=\sigma_{M}^{2}\,\beta_{i}\beta_{j}+\begin{cases}\sigma_{\varepsilon,i}^{2}&i=j\\0&i\neq j\end{cases}
  \]
  where $\sigma_{M}$ (1/yr) is the volatility of market growth and $\sigma_{\varepsilon,i}$ (1/yr) is the
  volatility of the residual growth of ticker $i$.
  {\setlength{\jot}{2pt}
  \begin{align*}
    &\min_{w}\ w^{\top}\Sigma\,w\quad\text{subject to}\quad\mathbf{1}^{\top}w=1 && \text{variance and budget}\\
    &2\,\Sigma\,w=\lambda\,\mathbf{1}\quad\Longrightarrow\quad w\propto\Sigma^{-1}\mathbf{1}
      && \text{first-order condition, multiplier }\lambda\\
    \Longrightarrow\quad&\boxed{w=\frac{\Sigma^{-1}\mathbf{1}}{\mathbf{1}^{\top}\Sigma^{-1}\mathbf{1}}}
  \end{align*}}%
  The code file then sets each negative weight to $0$ and divides the others by their sum, so that the
  portfolio holds no short positions.
\end{frame}

\begin{frame}{Under the Trading Limits}
  \vspace{0.3em}
  At each decision, the engine takes two steps, with its own covariance $\Sigma_{g}$:
  \begin{itemize}\setlength{\itemsep}{0.4em}
    \item \hilite{Mix}: minimize $w^{\top}\Sigma_{g}w$, over long-only weights that add to $s$, the largest share of the wealth that the other limits allow. This step leaves out the volatility limit.
    \item \hilite{Scale}: if the mix is more volatile than the limit $\sigma_{\max}$, scale it down, $w=c\,w_{\text{mix}}$ with $c=\min\big(1,\ \sigma_{\max}\big/\sqrt{w_{\text{mix}}^{\top}\Sigma_{g}w_{\text{mix}}}\big)$. The rest of the wealth stays in cash.
  \end{itemize}

  The engine's covariance also keeps the correlations between the residuals over the last 126 trading days:
  \[
    \Sigma_{g}=\sigma_{M}^{2}\,\beta\beta^{\top}+(1-\lambda)\,\hat{\Omega}_{\varepsilon}+\lambda\,\operatorname{diag}\big(\hat{\Omega}_{\varepsilon}\big)
  \]
  where $\hat{\Omega}_{\varepsilon}$ is the sample covariance of the residual growth rates, and $\lambda$
  between $0$ and $1$ shrinks it toward its diagonal.
\end{frame}

\begin{frame}{The Two Steps on a Made-Up Basket}
  \vspace{0.3em}
  Each point on the curve is a long-only portfolio of a made-up basket:
  \begin{center}
    \includegraphics{slides/figs/Fig-L15a-Deck-Two-Step-Minimum-Variance.pdf}
  \end{center}
  The minimum-variance mix is the leftmost point of the curve. This mix is within the volatility limit, and the
  engine holds the mix with $c=1$.
\end{frame}

\begin{frame}{Reading the Code}
  \vspace{0.3em}
  The rule's code file is
  \href{\repo/strategies/weighting/minimum_variance.jl}{\texttt{strategies/weighting/minimum\_variance.jl}}:
  \begin{itemize}\setlength{\itemsep}{0.5em}
    \item \texttt{weigh\_minimum\_variance(beta, sigma\_e, sigma\_m)} computes the weights with the single index covariance, before the trading limits. The engine does not call it.
    \item \texttt{MyL15aMinimumVarianceRule} is the type the engine calls. The rule has no settings of its own.
    \item \texttt{MPC.allocate(rule, p)} passes the decision's problem \texttt{p} to the engine's \texttt{MPC.MyMinVarRule}, which takes the two steps.
  \end{itemize}

  \vspace{0.5em}
  The \href{\repo/strategies/CHEME-5660-L15a-Strategy-Menu-Fall-2026.ipynb}{strategy menu} lists every rule, and the \href{\repo/engine/CHEME-5660-L15a-Trading-Engine-Fall-2026.ipynb}{trading engine} describes the shared limits.
\end{frame}

\end{document}
